\section{\methodname{} Details}
\label{app:iwr_details}
In this section, we provide some description to give a better understanding of how IWR addresses the challenge of mode switching. 


IWR does not change the CRL loss or its energy landscape directly. Instead, it changes the effective sampling distribution of contrastive updates. This provides an inductive bias toward mode-switching regions, giving more critic updates to interaction-relevant transitions rather than to easy passive motion or object-reaching behavior. In terms of our analysis (Section~\ref{sec: analysis}), IWR increases the probability that sampled positive futures cross the interaction bridge near the latent interaction time \(\tau_t\). Such samples expose the interaction shift directly, instead of observing it only after long passive propagation through powers of \(A_0\). Intuitively, based on the analysis of intrinsic reward \cite{bastankhah2026demystifying}, our paradigm shifts intrinsic reward-driven exploration from ``moving the object to the future'' to ``learning how to control the object to move to the future.'' This also connects to Proposition~\ref{prop:endpoint-energy-error} as increasing the coverage of interaction-centered samples reduces the error contribution from the interaction residual. Exact bridge samples provide direct supervision for the piecewise jump term, while the soft weighting in Eq.~\ref{eq:iwr} also covers nearby positives and reduces error accumulation from repeatedly propagating the interaction signal through \(A_0^k\).
Consequently, IWR reduces the passive bias of ordinary future sampling and improves the conditioning of interaction-relevant representation learning.


% We preserve the CRL loss landscape while adding an inductive bias that yields more updates on the interaction part, rather than on easy-to-learn object-reaching behavior or uncontrollable passive target-object motion.
% Intuitively, based on the analysis of intrinsic reward \cite{bastankhah2026demystifying}, our paradigm shifts intrinsic reward-driven exploration from ``moving the object to the future'' to ``learning how to control the object to move to the future.''
% Based on Proposition~\ref{prop:endpoint-energy-error}, a higher coverage of the interaction term would control the error bound $|\widehat{E}_k - E_k|$
% . On the one hand, exact interaction positives would help estimate the piecewise jump transfer by controlling the error bound of $\|e\|$.
% On the other hand, soft weighting label would sample positives close to the interaction point. Then we could reduce the accumulation of errors from the exponential iteration of $\|A_0^k\|$.
% As a result, interaction control would be more consistent and robust.